
\documentclass{article}
\usepackage{colortbl}
\usepackage{makecell}
\usepackage{multirow}
\usepackage{supertabular}

\begin{document}

\newcounter{utterance}

\centering \large Interaction Transcript for game `ifeval', experiment `ifeval', episode 161 with reflexion{-}llama{-}3.1{-}8b{-}instruct.
\vspace{24pt}

{ \footnotesize  \setcounter{utterance}{1}
\setlength{\tabcolsep}{0pt}
\begin{supertabular}{c@{$\;$}|p{.15\linewidth}@{}p{.15\linewidth}p{.15\linewidth}p{.15\linewidth}p{.15\linewidth}p{.15\linewidth}}
    \# & \multicolumn{2}{c}{Player} && \multicolumn{2}{c}{Game Master} \\
    \hline

    \theutterance \stepcounter{utterance}  
    & & & \multicolumn{4}{p{0.6\linewidth}}{
        \cellcolor[rgb]{0.9,0.9,0.9}{
            \makecell[{{p{\linewidth}}}]{
                \texttt{\tiny{[P1$\langle$GM]}}
                \texttt{You have attempted this same task before and did not succeed. These are the plans you wrote after those attempts.} \\
\texttt{Trial \#1: New plan: } \\
\\ 
\texttt{Upon reflection, I realize that my previous attempt was too focused on describing the performance itself, without considering the specific task of continuing and expanding the given sentence. I got caught up in the excitement of the performance and lost sight of the goal.} \\
\\ 
\texttt{To improve, I will take a step back and re{-}examine the original sentence: "A man is performing for a crowd on a stage, ..." My task is to continue and expand this sentence, not to create a new story or describe a different scene.} \\
\\ 
\texttt{Next time, I will focus on building on the existing sentence, using the given information as a starting point. I will ask myself questions like: What is the man doing on stage? What is the crowd's reaction? What is the atmosphere like?} \\
\\ 
\texttt{To achieve this, I will take a more structured approach, breaking down the task into smaller, manageable chunks. I will start by identifying the key elements of the original sentence, such as the man, the crowd, and the stage. Then, I will use these elements as a foundation to build on, adding new details and descriptions to continue and expand the sentence.} \\
\\ 
\texttt{For example, I might start by describing the man's actions on stage, such as singing or playing an instrument. I might then describe the crowd's reaction, including their emotions and physical responses. Finally, I might add more details about the stage and the atmosphere, using sensory language to bring the scene to life.} \\
\\ 
\texttt{By taking a more focused and structured approach, I hope to create a more cohesive and engaging continuation of the original sentence.} \\
\\ 
\\ 
\texttt{Can you please continue and expand: "A man is performing for a crowd on a stage, ..."} \\
\texttt{Use at least 60 sentences, but less than 600 words.} \\
            }
        }
    }
     \\ \\

    \theutterance \stepcounter{utterance}  
    & \multicolumn{4}{p{0.6\linewidth}}{
        \cellcolor[rgb]{0.561,0.737,0.561}{
            \makecell[{{p{\linewidth}}}]{
                \texttt{\tiny{[P1$\rangle$GM]}}
                \texttt{Upon re{-}examining the original sentence, I will focus on building on the existing information, using the given elements as a starting point. I will identify the key elements of the sentence, such as the man, the crowd, and the stage, and use these as a foundation to add new details and descriptions.} \\
\\ 
\texttt{The man, dressed in a black leather jacket and jeans, stood confidently at the microphone, his eyes scanning the crowd as he began to sing. His voice was deep and rich, filling the room with a sense of energy and excitement. The crowd, a mix of young and old, men and women, was entranced by his performance, their faces lit up with smiles and their hands clapping in time with the music.} \\
\\ 
\texttt{As the man sang, he moved across the stage, his feet tapping out the rhythm on the wooden floorboards. The stage itself was simple, with a single spotlight shining down on the man, but it was enough to create a sense of drama and spectacle. The crowd was packed tightly into the small venue, their bodies pressed up against the walls and the front of the stage.} \\
\\ 
\texttt{The air was thick with the smell of sweat and smoke, and the sound of the music was almost overwhelming, but the man's voice cut through it all, clear and strong. He was a natural performer, and he knew just how to work the crowd, drawing them in with his charisma and his talent.} \\
\\ 
\texttt{As he sang, the man's eyes met those of a young woman in the front row, her face alight with excitement. She was clapping along to the music, her hands moving in time with the beat, and the man caught her eye and smiled, his voice dipping down to a low, husky growl.} \\
\\ 
\texttt{The crowd was eating it up, their energy and enthusiasm feeding back into the man's performance, creating a sense of feedback and connection. He was feeding off their energy, using it to fuel his own performance, and it was creating a truly electric atmosphere.} \\
\\ 
\texttt{The stage was small, but the man made the most of it, using every inch of space to create a sense of drama and spectacle. He moved across the stage, his feet tapping out the rhythm, his voice soaring up to the rafters. The crowd was on its feet, cheering and clapping along to the music, and the man was feeding off their energy, using it to fuel his own performance.} \\
\\ 
\texttt{As the song came to an end, the man took a bow, his eyes scanning the crowd as he smiled and waved. The crowd erupted into applause, cheering and whistling for more, and the man grinned, his face alight with excitement. He knew he had them in the palm of his hand, and he was going to milk it for all it was worth.} \\
\\ 
\texttt{The next song was a slower, more introspective number, and the man's voice took on a softer, more melancholy tone. The crowd was entranced, their faces lit up with a sense of wonder and awe, and the man's eyes met those of the young woman in the front row, his voice dipping down to a low, husky growl.} \\
\\ 
\texttt{As the song came to an end, the man took a deep breath, his eyes scanning the crowd as he smiled and nodded. The crowd was on its feet, cheering and clapping along to the music, and the man was feeding off their energy, using it to fuel his own performance.} \\
\\ 
\texttt{The stage was small, but the man made the most of it, using every inch of space to create a sense of drama and spectacle. He moved across the stage, his feet tapping out the rhythm, his voice soaring up to the rafters. The crowd was on its feet, cheering and clapping along to the music, and the man was feeding off their energy, using it to fuel his own performance.} \\
\\ 
\texttt{As the night wore on, the man's performance only grew more intense and more passionate. He was a true performer, and he knew just how to work the crowd, drawing them in with his charisma and his talent. The crowd was entranced, their faces lit up with a sense of wonder and awe, and the man's eyes met those of the young woman in the front row, his voice dipping down to a low, husky growl.} \\
\\ 
\texttt{The music was a fusion of rock and blues, with a hint of country and folk. The man's voice was a powerful instrument, capable of conveying a wide range of emotions and moods. He sang of love and heartbreak, of joy and sorrow, and the crowd was entranced, their faces lit up with a sense of wonder and awe.} \\
\\ 
\texttt{As the night wore on, the man's performance only grew more intense and more passionate. He was a true performer, and he knew just how to work the crowd, drawing them in with his charisma and his talent. The crowd was on its feet, cheering and clapping along to the music, and the man was feeding off their energy, using it to fuel his own performance.} \\
\\ 
\texttt{The stage was small, but the man made the most of it, using every inch of space to create a sense of drama and spectacle. He moved across the stage, his feet tapping out the rhythm, his voice soaring up to the rafters. The crowd was entranced, their faces lit up with a sense of wonder and awe, and the man's eyes met those of the young woman in the front row, his voice dipping down to a low, husky growl.} \\
\\ 
\texttt{As the final song came to an end, the man took a deep bow, his eyes scanning the crowd as he smiled and nodded. The crowd erupted into applause, cheering and whistling for more, and the man grinned, his face alight with excitement. He knew he had given it his all, and he was proud of what he had accomplished.} \\
\\ 
\texttt{The crowd was on its feet, cheering and clapping along to the music, and the man was feeding off their energy, using it to fuel his own performance. He was a true performer, and he knew just how to work the crowd, drawing them in with his charisma and his talent. The music was a fusion of rock and blues, with a hint of country and folk, and the man's voice was a powerful instrument, capable of conveying a wide range of} \\
            }
        }
    }
    & & \\ \\

    \theutterance \stepcounter{utterance}  
    & & & \multicolumn{2}{p{0.3\linewidth}}{
        \cellcolor[rgb]{0.9,0.9,0.9}{
            \makecell[{{p{\linewidth}}}]{
                \texttt{\tiny{[GM$|$GM]}}
                {'length_constraints:number_sentences': {'relation': 'at least', 'num_sentences': 60}, 'length_constraints:number_words': {'relation': 'less than', 'num_words': 600}}
            }
        }
    }
    & & \\ \\

    \theutterance \stepcounter{utterance}  
    & & & \multicolumn{2}{p{0.3\linewidth}}{
        \cellcolor[rgb]{0.9,0.9,0.9}{
            \makecell[{{p{\linewidth}}}]{
                \texttt{\tiny{[GM$|$GM]}}
                \texttt{game\_result = LOSE} \\
            }
        }
    }
    & & \\ \\

    \theutterance \stepcounter{utterance}  
    & & \multicolumn{4}{p{0.6\linewidth}}{
        \cellcolor[rgb]{0.95,0.95,0.95}{
            \makecell[{{p{\linewidth}}}]{
                \texttt{\tiny{[P1$\leftrightarrow$P]}}
                \texttt{Post{-}episode Reflexion (not sent to GM):} \\
\texttt{New plan: } \\
\\ 
\texttt{Upon reflection, I realize that my previous attempt was too focused on describing the performance itself, without considering the specific task of continuing and expanding the given sentence. I got caught up in the excitement of the performance and lost sight of the goal.} \\
\\ 
\texttt{To improve, I will take a step back and re{-}examine the original sentence: "A man is performing for a crowd on a stage, ..." My task is to continue and expand this sentence, not to create a new story or describe a different scene.} \\
\\ 
\texttt{Next time, I will focus on building on the existing sentence, using the given information as a starting point. I will ask myself questions like: What is the man doing on stage? What is the crowd's reaction? What is the atmosphere like?} \\
\\ 
\texttt{To achieve this, I will take a more structured approach, breaking down the task into smaller, manageable chunks. I will start by identifying the key elements of the original sentence, such as the man, the crowd, and the stage. Then, I will use these elements as a foundation to build on, adding new details and descriptions to continue and expand the sentence.} \\
\\ 
\texttt{However, I also realize that my previous attempt was too repetitive and lacked variety. I repeated the same descriptions and ideas multiple times, without adding any new insights or perspectives.} \\
\\ 
\texttt{To address this, I will make a conscious effort to vary my descriptions and add more depth and nuance to the scene. I will use sensory language to bring the scene to life, incorporating details about the sights, sounds, and emotions of the performance.} \\
\\ 
\texttt{For example, I might describe the man's voice as "rich and velvety," or the crowd's reaction as "electric and frenzied." I might also add more details about the stage and the atmosphere, such as the smell of sweat and smoke, or the feeling of the crowd's energy.} \\
\\ 
\texttt{By taking a more focused and structured approach, and varying my descriptions to add more depth and nuance, I hope to create a more engaging and immersive continuation of the original sentence.} \\
            }
        }
    }
    & \\ \\

\end{supertabular}
}

\end{document}
